\subsection{Algebraic spectral retuning detects the collision order}
The integral-weight hypothesis can be removed if the comparison
retains the whole relative centralizer order. The actual
order-two contact then gives an obstruction that an abstract
special-fiber algebra alone does not detect.

\begin{finding}[\fid{F52}: the algebraic contact-order gate]\label{f:algebraiccontact}
Let $K/\C(v)$ be finite and let $F\in GL_N(K)$.
Fix $b>1$ and regular evaluations of this same algebraic
family such that, for every sufficiently large integer $n$,
\begin{equation}\label{eq:contactretuning}
 \operatorname{Spec}F_{b^n}
     =\operatorname{Spec}(F_b^n)
\end{equation}
as multisets. The evaluations may use different places.
No purity or integral spectral-weight hypothesis is imposed.

Let $R=\C[v,v^{-1}]_{(v-1)}$, let $R'$ be a discrete
valuation ring of $K$ above $R$, and suppose $F$ has
a regular invertible model on a finite free $R'$-module.
With
\[
 \delta=v+v^{-1},\quad D=\delta^2,\quad
 \Delta=\frac{(D-1)(D-4)}{D^2},
\]
its full relative centralizer cannot be isomorphic,
as an $R'$-algebra, to the actual order
\[
 A_{R'}=R'\times R'\times
 M_2\bigl(R'[x]/(x(x-\Delta))\bigr).
\]

There is a first-order algebraic model satisfying
\eqref{eq:contactretuning} whose abstract special
centralizer is nevertheless the same algebra as $A_1$.
It has nonintegral spectral weights and a different
relative order. The actual order embeds in its centralizer
and retains the actual rank-$64$ tensor module by restriction.
On that unchanged auxiliary module, adding ordinary linear
operators cannot make the embedded nonsaturated suborder
their simultaneous centralizer.
\end{finding}

\begin{proof}
\par\smallskip\noindent\textbf{Integral powers of the base eigenvalues.}\quad
Regard $F$ as a $\C(v)$-linear endomorphism of $K^N$.
Its characteristic polynomial over $\C(v)$ annihilates
the original $K$-linear matrix. Clear denominators to
obtain a nonzero $Q(v,Y)\in\C[v,Y]$ with $Q(v,F)=0$.
Outside finitely many parameter values its coefficients
are regular. For every base eigenvalue $\alpha$,
\eqref{eq:contactretuning} gives
\[
 Q(b^n,\alpha^n)=0
\]
eventually.

Expand $Q$ into distinct nonzero monomials
$q_{a,c}v^aY^c$. This is a finite power-sum relation
on roots $b^a\alpha^c$.
Vandermonde uniqueness forces every grouped coefficient
to vanish, so two distinct monomials have equal roots:
\[
 b^a\alpha^c=b^{a'}\alpha^{c'}.
\]
They cannot have $c=c'$, since then $b>1$ would force
$a=a'$ as well. Thus, after reversing signs if needed,
\[
 \alpha^L=b^a\quad\hbox{for some }L>0,\ a\in\Z.
\]
Choose one common $L$ for all base eigenvalues, including
their multiplicities, and write
\begin{equation}\label{eq:contactbasecharacters}
 \alpha_j^L=b^{a_j},\qquad a_j\in\Z.
\end{equation}
Their magnitudes are rational powers of $b$ and their
phases are roots of unity; this conclusion was proved
rather than assumed as a weight convention.

\par\smallskip\noindent\textbf{A monomial characteristic polynomial on a finite cover.}\quad
Adjoin $t$ with $v=t^L$.
The finite algebra
\[
 K\otimes_{\C(v)}\C(t)
\]
is a product of fields: in characteristic zero a
primitive-element polynomial for $K$ is separable,
and its distinct irreducible factors over $\C(t)$
give this product decomposition.
On at least one component $K''$ we claim
\begin{equation}\label{eq:contactcharpoly}
 \det(TI-F)=\prod_{j=1}^N(T-t^{a_j}).
\end{equation}

Indeed restrict \eqref{eq:contactretuning} to $n=Lr$,
and at $v=b^{Lr}$ prescribe the positive value $t=b^r$.
The evaluated spectrum is
\[
 \{\alpha_j^{Lr}\}_j=\{b^{a_jr}\}_j
                  =\{t^{a_j}\}_j.
\]
After discarding finitely many bad parameters, these
evaluations select factors of a finite model of the
tensor product. One component occurs for infinitely
many distinct $r$.
On that component each difference between the family
characteristic-polynomial coefficients and those
in \eqref{eq:contactcharpoly} has infinitely many zeros
over distinct $t=b^r$.

A nonzero algebraic function cannot have those zeros.
Its monic minimal polynomial over $\C(t)$ has nonzero
constant term; outside the finitely many coefficient
poles, a zero of the function forces a zero of that
nonzero rational constant term.
Thus every coefficient difference is zero, proving
\eqref{eq:contactcharpoly}.

The finite-model assertion needs no continuous sheet.
One may invert finitely many polynomials so that the
primitive-element integral model, discriminant, matrix
coefficients and the idempotents selecting the field
components are regular. The chosen evaluation on $K$
and $t=b^r$ then give a map from their tensor product
to $\C$, which selects one factor. Infinitely many
parameters survive every such finite deletion.
Other embeddings may multiply the resulting characters
by roots of unity; the local argument is unaffected.

\par\smallskip\noindent\textbf{The relative-center contradiction.}\quad
Extend $R'$ to a valuation ring $R''$ in that component
$K''$. The extension is a DVR and is torsion free,
hence flat, over $R'$.
The full centralizer is the kernel of
\[
 \operatorname{Mat}_N(R')\longrightarrow
 \operatorname{Mat}_N(R'),\qquad U\longmapsto UF-FU.
\]
Flat base change preserves this kernel.
Thus an assumed isomorphism with $A_{R'}$ would give
the same actual order over $R''$.
Put $e=\operatorname{ord}_{R''}(v-1)>0$.

The operator $F$ is central in its full centralizer.
Its component in the matrix-center factor must
therefore have the form $c_0+c_1x$, with
$c_0,c_1\in R''$.
On the two generic center branches its values are
\[
 \beta_1=c_0,\qquad
 \beta_2=c_0+c_1\Delta.
\]
Hence
\begin{equation}\label{eq:contactdivisibility}
 \beta_2-\beta_1\in\Delta R''.
\end{equation}
But
\[
 D-4=(v-v^{-1})^2
     =(v-1)^2(v+1)^2/v^2,
\]
and the remaining factor in $\Delta$ is a unit at
$v=1$. Therefore
\[
 \operatorname{ord}_{R''}(\Delta)=2e.
\]

The generic actual order is the split semisimple
product of two scalar and two two-by-two matrix
algebras. A central element in it is a tuple of
scalars, so its minimal polynomial is square free
after repeated values are removed. Thus $F$ is
diagonalizable generically.
Its values on the four primitive central blocks
must be distinct. Otherwise every linear map between
two nonzero eigenspaces with the same value would
commute with $F$, contradicting their status as
separate central blocks of the full centralizer.
In particular $\beta_1,\beta_2$ are distinct functions.

By \eqref{eq:contactcharpoly} they are two distinct
monomial characters $t^a,t^c$.
At the chosen point over $v=1$, $t=t_0$ with
$t_0^L=1$. The zero of $t^L-1$ at $t_0$ is simple,
so $\operatorname{ord}_{R''}(t-t_0)=e$.
The divisibility \eqref{eq:contactdivisibility}
forces $t_0^a=t_0^c$.
Distinctness gives $a\ne c$, and
\[
 \left.\frac{d}{dt}(t^a-t^c)\right|_{t=t_0}
       =(a-c)t_0^{a-1}\ne0.
\]
This holds also for negative exponents since
$t_0\ne0$. Thus
\[
 \operatorname{ord}_{R''}(\beta_2-\beta_1)=e,
\]
contradicting the required order at least $2e$.
This proves the full relative-order obstruction.

\par\smallskip\noindent\textbf{A sharp first-order model.}\quad
Let $v=t^3$ and choose a primitive cube root $\omega$
as the point $t=\omega$ over $v=1$. Put
\[
 \gamma_1=t^2,\qquad \gamma_2=t^5,\qquad
 \eta=\gamma_2-\gamma_1=t^2(v-1).
\]
Use the finite localization $T_0$ of $\C[t,t^{-1}]$
which inverts
\[
 \delta(v),\ h(v)=D(v)-1,\ D(v)-2,\quad
 t-1,\ t^2-1,\ t^5-1,\ t^2-t,\ t^5-t.
\]
All these factors are nonzero at $t=\omega$ and
at every positive real $t>1$.
This is an explicit finite coefficient base admitting
both the required local point and the arithmetic
evaluations.

Set
\[
 C_\eta=T_0[x]/(x(x-\eta)),\qquad
 J=C_\eta,\qquad F_J=\gamma_1+x.
\]
In the column basis $(1,x)$,
\[
 F_J=\begin{pmatrix}\gamma_1&0\\1&\gamma_2\end{pmatrix}.
\]
Its determinant is $t^7$, a unit.
Since $x=F_J-\gamma_1$, its centralizer on $J$ is
$C_\eta$, including at collision.
With $(T_0)_a$ denoting a scalar line of Frobenius $a$,
put
\[
 H=(T_0)_1\oplus(T_0)_t\oplus J^{\oplus2}.
\]
The sector differences are units:
the relevant determinants are
$(\gamma_1-1)(\gamma_2-1)$ and
$(\gamma_1-t)(\gamma_2-t)$.
Therefore
\begin{equation}\label{eq:firstordercentralizer}
 \End_F(H)=T_0\times T_0\times M_2(C_\eta).
\end{equation}

At $t=\omega$ the three eigenvalue sectors are
$1,\omega,\omega^2$, and $J$ remains a nontrivial
Jordan block in the last one.
The special algebra is abstractly
\[
 \C\times\C\times M_2(\C[\epsilon]/(\epsilon^2)),
\]
the same as $A_1$.
At $v=b>1$ choose $t_b=b^{1/3}$, and at $v=b^n$
choose $t_n=t_b^n$.
The spectrum consists of characters $1,t,t^2,t^5$
with multiplicities $1,1,2,2$, so it obeys full
spectral retuning.
The lower-left entry need not obey a fixed-basis
matrix-power identity; no such identity is claimed.

Under $v=Q^{1/2}$ these characters have spectral
weights $0,1/3,2/3,5/3$.
This is explicitly a nonintegral-weight model,
not a mixed model with integral Deligne weights.
The point $t=\omega$ is not the positive-real
branch limit $t=1$.
It does not contradict \fid{F51}.
The relative center in
\eqref{eq:firstordercentralizer} has first-order
contact, which explains why it also does not
contradict the present theorem.

\par\smallskip\noindent\textbf{The actual order as a specified subalgebra.}\quad
The element
\[
 c=\frac{\Delta}{\eta}
   =\frac{(D-1)(v-1)(v+1)^2}{D^2t^2v^2}
\]
belongs to $T_0$ and has order one at $t=\omega$.
There is an injection
\begin{equation}\label{eq:actualsuborder}
 T_0[x_{\rm act}]/(x_{\rm act}(x_{\rm act}-\Delta))
   \longrightarrow C_\eta,\qquad x_{\rm act}\longmapsto cx.
\end{equation}
Indeed $(cx)^2=\Delta(cx)$.
Its free-module matrix in the bases $(1,x_{\rm act})$
and $(1,x)$ is $\operatorname{diag}(1,c)$, proving
injectivity over the domain.
Keeping the two scalar sectors and taking two-by-two
matrices embeds the actual order into
\eqref{eq:firstordercentralizer}.

Locally at $t=\omega$, the center cokernel is
$R'/(c)$, of length one, and the matrix-sector
cokernel has length four.
The generic embedding is an isomorphism.
In the special fiber $x_{\rm act}$ maps to zero,
killing the actual four-dimensional matrix radical.
Thus the abstract special centralizer has the right
isomorphism type, but the actual relative embedding
is not that special isomorphism.

Let $L_0=C_\eta/(x)$ and $L_\eta=C_\eta/(x-\eta)$,
with Frobenii $\gamma_1,\gamma_2$, and define
\[
 K=(T_0)_1^{20}\oplus(T_0)_t^4
       \oplus L_0^{16}\oplus L_\eta^4.
\]
The cyclic source and unit sector differences give
\[
 \operatorname{Hom}_F(H,K)
  =\epsilon_+^{20}\oplus\epsilon_-^4
       \oplus\operatorname{Row}(L_0)^{16}
       \oplus\operatorname{Row}(L_\eta)^4,
\]
of rank $64$.
Under \eqref{eq:actualsuborder}, $x_{\rm act}$ acts
on these branches by $0,\Delta$.
This is the actual \fid{F21} tensor lattice after
the indicated coefficient extension, with its named
generators. The regular matrix comparison is defined
because $\delta,h,D-2$ are units on the chosen base.

Postcomposition on the multiplicity factors retains
the actual ordinary left operators as in \fid{F46}.
Over the domain, cross-branch intertwiners vanish
because $\eta$ is a nonzero divisor.
The multiplicity commutant is
$M_{20}(T_0)\times M_4(T_0)\times
 M_{16}(T_0)\times M_4(T_0)$, and its matrix units
lift to $K$.
At collision the branch multiplicities combine as
before. The spectra of $H$ and $K$ retune, with the
stated nonintegral weights, and their invariant Hom
has identity Frobenius.
This does not give the all-extension quantum-word
trace on that one fixed Hom space or the
\fid{F24} opposite-sign output.

\par\smallskip\noindent\textbf{Why extra linear commutators do not repair this embedding.}\quad
Localize the explicit model at $t=\omega$ and retain
this same free $H$. Let $B=\End_F(H)$ and let $A$
be the actual embedded suborder.
Any simultaneous centralizer of ordinary $R'$-linear
operators in $\End_{R'}(H)$ is saturated.
Indeed, if $d\ne0$ and $dU$ commutes with every
specified operator $G$, then
$d(UG-GU)=0$.
The matrix module is torsion free, so $UG=GU$ for
each $G$ and $U$ lies in that centralizer.

The embedded $A$ is not saturated.
Take $U=xE_{11}$ in the $M_2(C_\eta)$ sector.
It is a regular endomorphism in $B$ and is not in
$A$, since its $x$ coefficient is not divisible
by $c$. But $cU\in A$.
Therefore no additional ordinary linear operators
on this unchanged $H$ have simultaneous centralizer
equal to the displayed $A$.
Changing the module lattice or object, retaining
$A$ as explicit algebra data, or making a different
categorical comparison is outside this assertion.
\end{proof}

\paragraph{The exact remaining alternatives.}
The obstruction requires full spectral retuning of
one finite algebraic family and preservation of the
whole regular relative centralizer order.
It allows arbitrary evaluation sheets and auxiliary
rank and needs no purity assumption.
An abstract special-fiber isomorphism, a proper
specified subalgebra, a nonregular or nonflat model,
or a trace-only scalar comparison does not supply
those hypotheses.
The explicit model demonstrates this distinction
while retaining the actual tensor action by restriction.
Its saturation obstruction fixes that model and
embedding; it is not a prohibition on all extra
algebraic structures.
A regular-module wrapper with identity Frobenius
can tautologically retain algebra and module actions,
but does not produce the quantum-grading, canonical-basis
or path-sign comparison.

% Separate paragraph for Part III.
\paragraph{\fid{F52} in Grothendieck's language.}
The original order and its first-order overorder
define different coherent algebra sheaves on the
parameter trait. Their abstract special fibers agree,
while their embedded center lattices differ, with
the displayed nonzero cokernel lengths.
The affine module/sheaf framework is EGA~I,
Theorem~1.3.7, p.~86, Corollary~1.3.8, p.~87,
and Theorem~1.4.1, pp.~90--92.
The action comparison uses SGA~3, I, 4.7.2--3,
with cote~90, scan~17 as the audited framework anchor.
The finite coefficient extension is the type
allowed in VIII, \S2.2, p.~7.
The finite-cover, flat-kernel, valuation and
saturation arguments have been supplied explicitly
above; no unverified archive statement is substituted
for them.
The precise additional gate is preservation of the
whole relative order under the arithmetic comparison,
not merely its generic or special isomorphism type.

